%% kapa: exam=1 round=23 year=2012 period=1 form=A subject=civil_law subject_ko=민법 seq=7 number=7
%% source: Q-Net 감정평가사 기출문제 (cst00309) artlSeq=5052578 file=raw/1008/5052578/2012년 제23회 감정평가사 제1차시험 기출문제(1교시 A형).pdf
%% conversion: kordoc 4.24.51 -> tools/gen1.py (auto); review=proofread-sonnet; flags=none
\begin{question}{7}
부동산 취득시효에 관한 설명으로 옳지 않은 것은? (다툼이 있으면 판례에 의함)
\begin{choices}
  \item 중복등기 중 선등기가 원인무효가 아니어서 후등기가 무효로 된 경우, 말소되지 않은 후등기를 근거로 등기부취득시효의 완성을 주장할 수 없다.
  \item 시효기간 만료 전에 제3자가 등기명의를 넘겨받은 경우, 점유자는 시효기간완성 후에 그 제3자를 상대로 취득시효를 원인으로 소유권이전등기를 청구할 수 있다.
  \item 점유자가 취득시효를 주장하는 경우, 자주점유 여부에 관한 증명책임은 취득시효의 성립을 부정하는 자에게 있다.
  \item 자기 소유의 부동산은 취득시효의 대상이 될 수 없다.
  \item 점유자가 점유 개시 당시에 소유권 취득의 원인이 될 수 있는 법률행위가 없다는 사실을 알면서 타인 소유의 부동산을 무단점유한 경우, 취득시효는 인정되지 않는다.
\end{choices}
\end{question}
\answer{4}
